\part{Appendices}
\chapter{Simulation}
In this part i will write some useful information for running the simulation.
All test were done on debian based machine:
To see the graphical interface launched in the docker container you need to have to share the access to the X server using the command \texttt{xhost local:root}.


today the tool for converting the data are 
\href{https://www.ros.org/blog/noetic-eol/}{https://www.ros.org/blog/noetic-eol/}


implement also \href{https://github.com/ros-navigation/navigation2/tree/main/tools/planner\_benchmarking}{planning benchmark}

describe and \textcolor{red}{use} composition in nav2 to obtain better performance \cite{10132508}



Inoltre per gestire il traffico sulla rete utilizzeremo zenoh (Zero Overhead Network Protocol) ed il bridge per ros2 zenoh-bridge-ros2dds che è  pub/sub/query protocol unifying data in motion, data at rest and computations. It elegantly blends traditional pub/sub with geo distributed storage, queries and computations, while retaining a level of time and space efficiency that is well beyond any of the mainstream stacks. data liberator protocol. Zenoh liberates data in several dimensions. 


with chatgpt possiamo gestire il robot in due modi: si ferma ed esegue l'azione non communicando con rmf (la linea rimane bloccata dal suo passaggio) 
